\documentclass[11pt,a4paper]{article}

\usepackage[a4paper,margin=1in]{geometry}
\usepackage{fontspec}
\IfFontExistsTF{Times New Roman}{\setmainfont{Times New Roman}}{\setmainfont{TeX Gyre Termes}}
\IfFontExistsTF{Menlo}{\setmonofont{Menlo}}{\setmonofont{Latin Modern Mono}}
\usepackage{amsmath,amssymb}
\usepackage{booktabs}
\usepackage{longtable}
\usepackage{array}
\usepackage{enumitem}
\usepackage{fancyhdr}
\usepackage[round,authoryear]{natbib}
\usepackage[hidelinks]{hyperref}
\usepackage{xurl}
\usepackage{microtype}

\setlength{\parindent}{0pt}
\setlength{\parskip}{0.65em}
\setlist[itemize]{leftmargin=1.4em}
\newcommand{\file}[1]{\path{#1}}
\newcommand{\code}[1]{\path{#1}}
\newcolumntype{L}[1]{>{\raggedright\arraybackslash}p{#1}}

\pagestyle{fancy}
\fancyhf{}
\lhead{\small Surrogate Dataset Design}
\rhead{\small Planning Only}
\cfoot{\thepage}
\renewcommand{\headrulewidth}{0.4pt}

\title{\textbf{Surrogate Dataset Design Plan}\\
\large Planning Gate Before Any Dataset Construction}
\author{Codex-assisted surrogate planning report}
\date{June 17, 2026}

\begin{document}
\maketitle

\begin{abstract}
This report defines the planned schema, parameter grid, label policy, mask policy, diagnostics,
split strategy, and construction gate for a future surrogate dataset. It approves dataset design
only. It does not generate data, labels, arrays, neural-network files, training tables, broad
parameter sweeps, or surrogate models. All future labels must come from the validated baseline
American Crank--Nicolson / projected SOR solver unless a later human-review gate approves a
separately labelled alternative.
\end{abstract}

\tableofcontents
\newpage

\section{Purpose of the Dataset Design Plan}

The application-strength gate concluded \code{READY_FOR_DATASET_PLANNING_WITH_LIMITATIONS}. The
next step is therefore a written dataset design plan, not dataset construction. This report
records the planned future schema and quality policy so that a later construction ticket can be
reviewed before it creates any arrays or labels.

The plan follows the project's evidence ladder: classical American option validation, pilot
stress-map diagnostics, quantitative risk analytics, and then dataset planning. This order is
important because a neural surrogate trained on unreliable numerical labels would only learn
unreliable structure. The mathematical background remains the Black--Scholes American
free-boundary and finite-difference setting
\citep{black_scholes_1973,merton_1973,brennan_schwartz_1977,wilmott_howison_dewynne_1995}.

\section{Link to Solver Validation and Application-Strength Gate}

Ticket 12 passed solver validation with limitations. The downstream design gate approved only a
small pilot. Pilot 01 and quantitative risk analytics then showed coherent application messages:
put boundaries moved lower as volatility increased, dividend-call boundaries moved lower as
dividend yield increased, high volatility increased PSOR effort, and LCP diagnostics remained
small in the tested regimes.

Those results support dataset planning. They do not support immediate generation, broad sweeps,
neural training, production Greek labels, or final paper claims.

\begin{longtable}{L{0.30\textwidth} L{0.48\textwidth}}
\toprule
Prior gate & Constraint carried into this plan \\
\midrule
\endhead
Ticket 12 synthesis & Solver is validated with limitations, not production certified. \\
Application-strength gate & Planning is allowed; construction and training remain blocked. \\
Pilot 01 & Use baseline CN/PSOR, \(K=1\), and full diagnostics. \\
Quantitative analytics & Preserve boundary, Gamma, PSOR, runtime, and LCP metadata. \\
Reference integration & Use verified citations and avoid unsupported literature claims. \\
\bottomrule
\end{longtable}

\section{Dataset Purpose and Non-Generation Statement}

The future dataset is intended to support constrained surrogate experiments for American option
price and continuation-premium surfaces, boundary diagnostics, and possibly masked Greek
diagnostics. It should help answer whether a surrogate can accelerate scenario evaluation while
respecting the payoff obstacle, continuation premium, and boundary-aware risk structure.

This step does not produce the dataset. The files in \file{reports/05_surrogate/} are planning
artifacts only. No \code{.npz}, \code{.npy}, \code{.pt}, training table, model file, or generated
label file is created by this design plan.

\section{Solver Source and Default Numerical Settings}

All future labels must come from the validated baseline solver:
\begin{center}
\code{solver_name = american_crank_nicolson_psor_price}\\
\code{solver_variant = baseline_cn_psor}
\end{center}

Rannacher-smoothed outputs are not the default. Ticket 10A kept Rannacher as an optional
comparison tool only \citep{rannacher_1984}. If Rannacher is ever included in a later dataset, it
must be separately labelled and must not be pooled silently with baseline CN/PSOR labels.

The first construction plan should inherit the downstream defaults unless human review changes
them: normalized strike \(K=1\), default \(S_{\max}=4K\), default \(M=N=120\), and selected
higher-grid confirmation at \(M=N=180\) before broad generation.

\section{Option Families}

The planned dataset must include:
\begin{itemize}
  \item American puts;
  \item dividend-paying American calls.
\end{itemize}

Both families are required because the application analytics produced meaningful but different
messages for puts and dividend-paying calls. A pooled dataset that hides dividend-call boundary
weakness would not serve the project. No-dividend American calls remain useful as validation
controls, but they are not the core first dataset family.

\section{Feature Schema}

The primary learning features are normalized variables:
\[
x = \left(\log(S/K),\ \tau/T,\ r,\ q,\ \sigma,\ T,\ \text{is\_call}\right).
\]
This matches the methodology handout's normalized feature recommendation while keeping enough
metadata for solver audit.

\begin{longtable}{L{0.24\textwidth} L{0.33\textwidth} L{0.16\textwidth} L{0.17\textwidth}}
\toprule
Feature & Definition & Required & Notes \\
\midrule
\endhead
\code{log_moneyness} & \(\log(S/K)\). & yes & Primary spot feature. \\
\code{tau_fraction} & \(\tau/T\). & yes & Normalized time-to-maturity. \\
\code{r} & Risk-free rate. & yes & Scale consistently in preprocessing. \\
\code{q} & Continuous dividend yield. & yes & Needed for both option families. \\
\code{sigma} & Volatility. & yes & Stress-sensitive regime feature. \\
\code{T} & Maturity in years. & yes & Planned grid includes four values. \\
\code{is_call} & 1 for call, 0 for put. & yes & Numeric option-family indicator. \\
\code{option_type} & \code{put} or \code{call}. & yes & Human-readable audit field. \\
\code{S_over_K} & \(S/K\). & optional & Redundant but useful for reporting. \\
\code{Smax}, \code{M}, \code{N} & Solver domain and grid metadata. & yes & Required for reproducibility. \\
\bottomrule
\end{longtable}

The companion CSV version is listed in the reproducibility and source record.

\section{Label Schema}

The first design separates reliable price-style labels from diagnostic labels. This protects the
project from treating every computed quantity as equally trustworthy.

\begin{longtable}{L{0.30\textwidth} L{0.25\textwidth} L{0.16\textwidth} L{0.19\textwidth}}
\toprule
Label & Definition & Reliability & Allowed use \\
\midrule
\endhead
\code{value_over_K} & \(U/K\). & core reliable & Primary future label. \\
\code{payoff_over_K} & \(\Phi/K\). & analytic reliable & Reference label. \\
\code{premium_over_K} & \((U-\Phi)/K\). & reliable with caution & Primary future label. \\
\code{exercise_indicator} & Threshold-derived exercise flag. & diagnostic & Metadata only at first. \\
\code{boundary_spot_over_K} & Extracted boundary divided by \(K\). & diagnostic & Boundary metadata. \\
\code{delta} & Finite-difference Delta. & cautioned & Optional diagnostic label. \\
\code{scaled_gamma} & \(K\Gamma\). & high caution & Masked diagnostic only. \\
\bottomrule
\end{longtable}

Price and premium labels are the most reliable first targets. Boundary labels remain
threshold-based. Gamma labels are fragile near payoff kinks, maturity, and extracted boundaries;
they must never be used unmasked. The companion CSV version is listed in the reproducibility and
source record.

\section{Parameter Grid Plan}

The first planned grid is intentionally small:
\[
T \in \{0.25,0.5,1.0,2.0\},\quad
\sigma \in \{0.20,0.40,0.60\},\quad
r \in \{0.01,0.05,0.10\}.
\]
For dividend yield:
\[
q \in \{0.00,0.03,0.06,0.10\}.
\]
For dividend-paying calls, the first construction plan should ensure positive-dividend regimes
are included. The \(q=0\) no-dividend call is mainly a validation control and should not dominate
the dividend-call dataset.

This grid is a plan, not generated data. The companion CSV version is listed in the reproducibility
and source record.

\section{Spot-Time Sampling Plan}

Future solves should keep the full solver domain for boundary conditions and diagnostics. The
main learning and reporting region should be:
\[
S/K \in [0.4,1.8],\qquad \tau/T \in [0,1].
\]

The recommended first version is a regular-grid design: retain grid nodes in the target region
because they are reproducible and easy to audit. Boundary-near, kink-near, and maturity rows
should be marked by masks rather than oversampled in this design step. Random samples from
surfaces and boundary-focused oversampling may be considered later, but only after the first
construction plan fixes masks, splits, and diagnostics.

\section{Mask Policy}

Masks are mandatory metadata. They are not cosmetic fields. They tell future surrogate code and
human reviewers which labels are safe to use for which purpose.

\begin{longtable}{L{0.31\textwidth} L{0.36\textwidth} L{0.21\textwidth}}
\toprule
Mask & Meaning & Required use \\
\midrule
\endhead
\code{payoff_kink_near} & Node is near \(S=K\), where payoff slope changes. & Greek caution. \\
\code{boundary_near} & Node is near the extracted free-boundary diagnostic. & Greek and boundary caution. \\
\code{maturity_row} & Node belongs to \(\tau=0\). & Greek caution. \\
\code{strict_interior} & Node avoids endpoints, kink, maturity, and boundary bands. & Safer Greek summaries. \\
\code{gamma_allowed_mask} & Node is allowed for Gamma target use. & Required for Gamma. \\
\code{delta_allowed_mask} & Node is allowed for Delta target use. & Required for Delta. \\
\code{exercise_region} & Threshold-derived exercise-like region. & Boundary metadata. \\
\code{continuation_region} & Threshold-derived continuation-like region. & Boundary metadata. \\
\bottomrule
\end{longtable}

The mask policy follows the Ticket 09 boundary extraction and Ticket 10 Greek-diagnostic logic.
It also follows the handout warning not to train heavily on Gamma exactly at payoff kinks or
extracted boundaries.

\section{Diagnostic Metadata Policy}

Every future solved regime must store diagnostics before any label is accepted:
\begin{itemize}
  \item all PSOR steps converged;
  \item maximum and mean PSOR iterations;
  \item maximum final PSOR update;
  \item maximum obstacle violation;
  \item maximum equation violation;
  \item maximum absolute complementarity product;
  \item boundary found count and boundary threshold;
  \item \(M\), \(N\), \(S_{\max}\), \(\Delta S\), and \(\Delta\tau\);
  \item solver name and solver variant;
  \item downstream use status.
\end{itemize}

If a regime fails convergence or diagnostic checks, it must be excluded from label export until
human review decides how to handle it. The companion CSV version is listed in the reproducibility
and source record.

\section{Train/Test Split Policy}

The main split must be regime-level: hold out entire parameter combinations. A random point split
inside a solved surface is allowed only as a secondary debugging diagnostic. Neighboring grid
points from the same surface are highly correlated; using them in both train and test can make a
surrogate look better than it really is.

\begin{longtable}{L{0.26\textwidth} L{0.37\textwidth} L{0.25\textwidth}}
\toprule
Split type & Purpose & Priority \\
\midrule
\endhead
Regime holdout & Test generalization to unseen parameter combinations. & primary \\
Option-family balance & Ensure puts and dividend calls are both represented. & primary \\
Stress-regime holdout & Test high-volatility, high-dividend, or low-rate regions. & primary \\
Random point split & Debug interpolation within solved surfaces. & secondary only \\
Boundary-band holdout & Optional future boundary-specific diagnostic. & review required \\
\bottomrule
\end{longtable}

The machine-readable version is \file{reports/05_surrogate/dataset_split_policy.csv}.

\section{Storage Recommendation}

The future construction ticket may use compressed arrays such as \code{.npz} plus CSV or JSON
manifests. A reasonable future package would include one array file for numeric samples and one
human-readable manifest for regimes, diagnostics, masks, and split membership.

This report does not create that storage. No array file, model file, training table, or label
export is produced here.

\section{Dataset Construction Gate}

Dataset construction may start only after human review approves:
\begin{itemize}
  \item parameter grid;
  \item feature schema;
  \item label schema;
  \item mask policy;
  \item diagnostic policy;
  \item split strategy;
  \item storage format;
  \item handling of failed diagnostics.
\end{itemize}

Approval of this design plan is not approval to train models. A later surrogate plan must review
model architecture, losses, metrics, held-out regime tests, and PSOR spot-check policy.

\section{Blocked Actions}

The following remain blocked:
\begin{itemize}
  \item actual dataset construction;
  \item \code{.npz}, \code{.npy}, \code{.pt}, or training-table export;
  \item neural surrogate training;
  \item broad production-scale sweeps;
  \item unmasked Gamma labels;
  \item final paper claims based on a dataset that does not yet exist.
\end{itemize}

\section{Limitations}

The planned grid is deliberately small. It should be treated as a first design, not a production
coverage claim. Boundary labels are threshold-based. Gamma labels remain fragile. Runtime and
PSOR effort may change under larger grids or wider parameter ranges. The plan also does not yet
specify exact train/test regime IDs; that should be done in the construction plan before any data
is generated.

\section{Recommended Next Step}

The recommended next step is human review of the five schema and policy CSVs plus this report. If
accepted, the next technical artifact should be a dataset construction plan that names the exact
regimes, storage format, diagnostic failure rules, and split assignments. Only after that plan is
approved should the project generate any dataset files.

\section{Reproducibility and Source Record}

\begin{longtable}{L{0.34\textwidth} L{0.50\textwidth}}
\toprule
Item & Record \\
\midrule
\endhead
Feature schema & \file{reports/05_surrogate/dataset_feature_schema.csv}. \\
Label schema & \file{reports/05_surrogate/dataset_label_schema.csv}. \\
Parameter grid plan & \file{reports/05_surrogate/dataset_parameter_grid_plan.csv}. \\
Split policy & \file{reports/05_surrogate/dataset_split_policy.csv}. \\
Diagnostic policy & \file{reports/05_surrogate/dataset_diagnostic_policy.csv}. \\
Backlog & \file{reports/05_surrogate/dataset_generation_backlog.md}. \\
Design report & \file{reports/05_surrogate/surrogate_dataset_design_plan.tex}. \\
Shared bibliography & \file{reports/03_solver/references.bib}. \\
\bottomrule
\end{longtable}

\bibliographystyle{plainnat}
\bibliography{reports/03_solver/references}

\end{document}
